% Documento de entrega — infografía interactiva sobre modelos de lenguaje.
\documentclass[a4paper,10pt]{article}

\usepackage[spanish]{babel}
\usepackage[T1]{fontenc}
\usepackage{lmodern}
\usepackage{microtype}
\usepackage[a4paper,top=2.4cm,bottom=2.4cm,left=2.4cm,right=2.4cm]{geometry}
\usepackage{xcolor}
\usepackage{url}
\usepackage[pdfauthor={Jheison Martinez Bolivar},pdftitle={Documento de entrega — Infografía interactiva sobre modelos de lenguaje}]{hyperref}
\usepackage[compact]{titlesec}
\usepackage{parskip}
\setlength{\parindent}{0pt}

\newcommand{\linkcolor}{blue!55!black}
\hypersetup{urlcolor=blue!55!black}

\begin{document}
% ─── Ficha de portada ───────────────────────────────────────────────────────
\begin{center}
\setlength{\fboxsep}{14pt}
\fcolorbox{gray}{gray!6}{%
\begin{minipage}{0.82\linewidth}
\centering
{\footnotesize\scshape documento de entrega}\\[6pt]
{\Large\bfseries Infografía interactiva}\\[2pt]
{\large modelos de lenguaje, 2017--2026}\\[10pt]
{\bfseries Jheison Martinez Bolivar}\\[2pt]
{\footnotesize octubre de 2026}\\[10pt]
{\large\href{https://jheisonmb.github.io/modelos-de-lenguaje/}{jheisonmb.github.io/modelos-de-lenguaje}}
\end{minipage}}
\end{center}

\vspace{1.5em}

\section*{El enlace}
La infografía está publicada como GitHub Pages en
\href{https://jheisonmb.github.io/modelos-de-lenguaje/}{jheisonmb.github.io/modelos-de-lenguaje};
el código fuente de la infografía está en el repositorio \url{https://github.com/JheisonMB/modelos-de-lenguaje}.

\section*{Justificación de las herramientas elegidas}

\subsection*{GitHub Pages}
Uno de los criterios era que la entrega fuera explorable, no estática: filtros,
fichas, tooltips, simulador de muestreo y balanza de evidencia. Un PDF puede
mostrar el diseño, pero no responder al cursor ni al foco del teclado.

\medskip

GitHub Pages permite publicar un sitio estático sin backend ni build, directamente
desde el repositorio; la infografía misma es open source: su código vive ahí.

\subsection*{\href{https://univerlab.org/texforge}{texforge}}
Para este documento se usó texforge, un compilador LaTeX independiente que compila
a PDF en un comando, con linter y corrector ortográfico en español incluidos. El resultado es un PDF con fuentes embebidas
que el receptor abre sin instalar nada.

\end{document}
